\documentclass[10pt,a4paper]{amsart}
\usepackage{amsmath,amssymb,amsthm,graphicx}
\usepackage{verbatim}
\newenvironment{reference}{\comment}{\endcomment}
%\newenvironment{reference}{}{}
\newenvironment{slogan}{\comment}{\endcomment}
\newenvironment{history}{\comment}{\endcomment}

% For commutative diagrams we use Xy-pic
\usepackage[all]{xy}

% We use 2cell for 2-commutative diagrams.
\xyoption{2cell}
% Theorem environments.
%
\theoremstyle{plain}
\newtheorem{theorem}[subsection]{Theorem}
\newtheorem{proposition}[subsection]{Proposition}
\newtheorem{lemma}[subsection]{Lemma}

\theoremstyle{definition}
\newtheorem{definition}[subsection]{Definition}
\newtheorem{example}[subsection]{Example}
\newtheorem{exercise}[subsection]{Exercise}
\newtheorem{situation}[subsection]{Situation}

\theoremstyle{remark}
\newtheorem{remark}[subsection]{Remark}
\newtheorem{remarks}[subsection]{Remarks}

\numberwithin{equation}{subsection}

% Macros
%
\def\lim{\mathop{\mathrm{lim}}\nolimits}
\def\colim{\mathop{\mathrm{colim}}\nolimits}
\def\Spec{\mathop{\mathrm{Spec}}}
\def\Hom{\mathop{\mathrm{Hom}}\nolimits}
\def\Ext{\mathop{\mathrm{Ext}}\nolimits}
\def\SheafHom{\mathop{\mathcal{H}\!\mathit{om}}\nolimits}
\def\SheafExt{\mathop{\mathcal{E}\!\mathit{xt}}\nolimits}
\def\Sch{\mathit{Sch}}
\def\Mor{\mathop{\mathrm{Mor}}\nolimits}
\def\Ob{\mathop{\mathrm{Ob}}\nolimits}
\def\Sh{\mathop{\mathit{Sh}}\nolimits}
\def\NL{\mathop{N\!L}\nolimits}
\def\CH{\mathop{\mathrm{CH}}\nolimits}
\def\proetale{{pro\text{-}\acute{e}tale}}
\def\etale{{\acute{e}tale}}
\def\QCoh{\mathit{QCoh}}
\def\Ker{\mathop{\mathrm{Ker}}}
\def\Im{\mathop{\mathrm{Im}}}
\def\Coker{\mathop{\mathrm{Coker}}}
\def\Coim{\mathop{\mathrm{Coim}}}

% Boxtimes
%
\DeclareMathSymbol{\boxtimes}{\mathbin}{AMSa}{"02}

%
% Macros for moduli stacks/spaces
%
\def\QCohstack{\mathcal{QC}\!\mathit{oh}}
\def\Cohstack{\mathcal{C}\!\mathit{oh}}
\def\Spacesstack{\mathcal{S}\!\mathit{paces}}
\def\Quotfunctor{\mathrm{Quot}}
\def\Hilbfunctor{\mathrm{Hilb}}
\def\Curvesstack{\mathcal{C}\!\mathit{urves}}
\def\Polarizedstack{\mathcal{P}\!\mathit{olarized}}
\def\Complexesstack{\mathcal{C}\!\mathit{omplexes}}
% \Pic is the operator that assigns to X its picard group, usage \Pic(X)
% \Picardstack_{X/B} denotes the Picard stack of X over B
% \Picardfunctor_{X/B} denotes the Picard functor of X over B
\def\Pic{\mathop{\mathrm{Pic}}\nolimits}
\def\Picardstack{\mathcal{P}\!\mathit{ic}}
\def\Picardfunctor{\mathrm{Pic}}
\def\Deformationcategory{\mathcal{D}\!\mathit{ef}}
\usepackage{cleveref}
\crefname{originalreference}{original reference}{original references}
\Crefname{originalreference}{Original reference}{Original references}
\makeatletter
\newlabel{stacks-morphisms-lemma-composition-smooth}{{075V}{source}{}{sourcepointer-1}{}}
\newlabel{stacks-morphisms-lemma-composition-smooth@cref}{{[originalreference][0][]075V}{[1][0][]source}{}{}{}}
\makeatother
\crefname{originalreference}{original reference}{original references}
\Crefname{originalreference}{Original reference}{Original references}
\makeatletter
\newlabel{sites-modules-lemma-local-pullback}{{03DO}{source}{}{sourcepointer-2}{}}
\newlabel{sites-modules-lemma-local-pullback@cref}{{[originalreference][0][]03DO}{[1][0][]source}{}{}{}}
\makeatother
\crefname{originalreference}{original reference}{original references}
\Crefname{originalreference}{Original reference}{Original references}
\makeatletter
\newlabel{stacks-sheaves-lemma-surjective-flat-locally-finite-presentation}{{06XH}{source}{}{sourcepointer-3}{}}
\newlabel{stacks-sheaves-lemma-surjective-flat-locally-finite-presentation@cref}{{[originalreference][0][]06XH}{[1][0][]source}{}{}{}}
\makeatother
\crefname{originalreference}{original reference}{original references}
\Crefname{originalreference}{Original reference}{Original references}
\makeatletter
\newlabel{sites-modules-lemma-local-final-object}{{03DN}{source}{}{sourcepointer-4}{}}
\newlabel{sites-modules-lemma-local-final-object@cref}{{[originalreference][0][]03DN}{[1][0][]source}{}{}{}}
\makeatother
\crefname{originalreference}{original reference}{original references}
\Crefname{originalreference}{Original reference}{Original references}
\makeatletter
\newlabel{stacks-sheaves-lemma-characterize-quasi-coherent-bis}{{0EM8}{source}{}{sourcepointer-5}{}}
\newlabel{stacks-sheaves-lemma-characterize-quasi-coherent-bis@cref}{{[originalreference][0][]0EM8}{[1][0][]source}{}{}{}}
\makeatother
\crefname{originalreference}{original reference}{original references}
\Crefname{originalreference}{Original reference}{Original references}
\makeatletter
\newlabel{stacks-sheaves-lemma-characterize-quasi-coherent}{{06WI}{source}{}{sourcepointer-6}{}}
\newlabel{stacks-sheaves-lemma-characterize-quasi-coherent@cref}{{[originalreference][0][]06WI}{[1][0][]source}{}{}{}}
\makeatother
\crefname{originalreference}{original reference}{original references}
\Crefname{originalreference}{Original reference}{Original references}
\makeatletter
\newlabel{stacks-morphisms-lemma-composition-flat}{{06PX}{source}{}{sourcepointer-7}{}}
\newlabel{stacks-morphisms-lemma-composition-flat@cref}{{[originalreference][0][]06PX}{[1][0][]source}{}{}{}}
\makeatother
\crefname{originalreference}{original reference}{original references}
\Crefname{originalreference}{Original reference}{Original references}
\makeatletter
\newlabel{stacks-morphisms-lemma-composition-finite-presentation}{{06Q3}{source}{}{sourcepointer-8}{}}
\newlabel{stacks-morphisms-lemma-composition-finite-presentation@cref}{{[originalreference][0][]06Q3}{[1][0][]source}{}{}{}}
\makeatother
\crefname{originalreference}{original reference}{original references}
\Crefname{originalreference}{Original reference}{Original references}
\makeatletter
\newlabel{stacks-sheaves-lemma-pullback-quasi-coherent}{{06WH}{source}{}{sourcepointer-9}{}}
\newlabel{stacks-sheaves-lemma-pullback-quasi-coherent@cref}{{[originalreference][0][]06WH}{[1][0][]source}{}{}{}}
\makeatother
\title{Quasi-coherence descends along smooth (resp. flat lfp) coverings}
\author{The Stacks Project Authors}
\date{stacks-project:stacks-cohomology.tex:local-snapshot-20261012; GNU FDL 1.2 or later}
\begin{document}
\maketitle
\noindent\textit{Editorial note.} All delivered mathematical text is the author's own, copied verbatim from the pinned Stacks Project chapter file stacks-cohomology.tex; no statement or proof step was rewritten, reordered or removed.\par\medskip
\noindent\textit{Editorial note.} The counted span cites the lemmas recorded in assembly.external\_references with canonical tags as navigation pointers.\par\medskip
\noindent\textit{Allowed standard prerequisites (editorial).}\begin{itemize}
\item Lifting over surjective flat lfp morphisms. Sheaves on Stacks, Lemma stacks-sheaves-lemma-surjective-flat-locally-finite-presentation. Provides the covering with lifts to U.
\item Local final object detection. Modules on Sites, Lemma sites-modules-lemma-local-final-object. qc on a covering implies qc downstairs.
\item Characterization of quasi-coherent modules on stacks. Sheaves on Stacks, Lemmas stacks-sheaves-lemma-characterize-quasi-coherent-bis and -quasi-coherent. Concludes qc from the pointwise check.
\end{itemize}\medskip
\ifdefined\hypertarget\hypertarget{sourcepointer-1}{}\fi
\noindent\textit{Source reference.} Original reference 075V: sources/stacks-chow/tags.txt, source lines 9144--9144. Cross-chapter standard prerequisite cited by the author inside the counted span; it is used as a pointer and its own proof is not part of this item.\par\medskip
\ifdefined\hypertarget\hypertarget{sourcepointer-2}{}\fi
\noindent\textit{Source reference.} Original reference 03DO: sources/stacks-chow/tags.txt, source lines 4427--4427. Cross-chapter standard prerequisite cited by the author inside the counted span; it is used as a pointer and its own proof is not part of this item.\par\medskip
\ifdefined\hypertarget\hypertarget{sourcepointer-3}{}\fi
\noindent\textit{Source reference.} Original reference 06XH: sources/stacks-chow/tags.txt, source lines 8851--8851. Cross-chapter standard prerequisite cited by the author inside the counted span; it is used as a pointer and its own proof is not part of this item.\par\medskip
\ifdefined\hypertarget\hypertarget{sourcepointer-4}{}\fi
\noindent\textit{Source reference.} Original reference 03DN: sources/stacks-chow/tags.txt, source lines 4426--4426. Cross-chapter standard prerequisite cited by the author inside the counted span; it is used as a pointer and its own proof is not part of this item.\par\medskip
\ifdefined\hypertarget\hypertarget{sourcepointer-5}{}\fi
\noindent\textit{Source reference.} Original reference 0EM8: sources/stacks-chow/tags.txt, source lines 18298--18298. Cross-chapter standard prerequisite cited by the author inside the counted span; it is used as a pointer and its own proof is not part of this item.\par\medskip
\ifdefined\hypertarget\hypertarget{sourcepointer-6}{}\fi
\noindent\textit{Source reference.} Original reference 06WI: sources/stacks-chow/tags.txt, source lines 8817--8817. Cross-chapter standard prerequisite cited by the author inside the counted span; it is used as a pointer and its own proof is not part of this item.\par\medskip
\ifdefined\hypertarget\hypertarget{sourcepointer-7}{}\fi
\noindent\textit{Source reference.} Original reference 06PX: sources/stacks-chow/tags.txt, source lines 8586--8586. Cross-chapter standard prerequisite cited by the author inside the counted span; it is used as a pointer and its own proof is not part of this item.\par\medskip
\ifdefined\hypertarget\hypertarget{sourcepointer-8}{}\fi
\noindent\textit{Source reference.} Original reference 06Q3: sources/stacks-chow/tags.txt, source lines 8592--8592. Cross-chapter standard prerequisite cited by the author inside the counted span; it is used as a pointer and its own proof is not part of this item.\par\medskip
\ifdefined\hypertarget\hypertarget{sourcepointer-9}{}\fi
\noindent\textit{Source reference.} Original reference 06WH: sources/stacks-chow/tags.txt, source lines 8816--8816. Cross-chapter standard prerequisite cited by the author inside the counted span; it is used as a pointer and its own proof is not part of this item.\par\medskip
\begin{lemma}
\label{lemma-check-qc-on-etale-covering}
Let $\mathcal{X}$ be an algebraic stack.
\begin{enumerate}
\item Let $f_j : \mathcal{X}_j \to \mathcal{X}$ be a family of smooth
morphisms of algebraic stacks with
$|\mathcal{X}| =\bigcup |f_j|(|\mathcal{X}_j|)$.
Let $\mathcal{F}$ be a sheaf of $\mathcal{O}_\mathcal{X}$-modules
on $\mathcal{X}_\etale$. If each $f_j^{-1}\mathcal{F}$
is quasi-coherent, then so is $\mathcal{F}$.
\item Let $f_j : \mathcal{X}_j \to \mathcal{X}$ be a family of flat and
locally finitely presented morphisms of algebraic stacks with
$|\mathcal{X}| =\bigcup |f_j|(|\mathcal{X}_j|)$.
Let $\mathcal{F}$ be a sheaf of $\mathcal{O}_\mathcal{X}$-modules
on $\mathcal{X}_{fppf}$. If each $f_j^{-1}\mathcal{F}$
is quasi-coherent, then so is $\mathcal{F}$.
\end{enumerate}
\end{lemma}

\begin{proof}
Proof of (1). We may replace each of the algebraic stacks $\mathcal{X}_j$
by a scheme $U_j$ (using that any algebraic stack has a smooth covering by
a scheme and that compositions of smooth morphisms are smooth, see
Morphisms of Stacks, Lemma \ref{stacks-morphisms-lemma-composition-smooth}).
The pullback of $\mathcal{F}$ to $(\Sch/U_j)_\etale$ is still
quasi-coherent, see
Modules on Sites, Lemma \ref{sites-modules-lemma-local-pullback}.
Then $f = \coprod f_j : U = \coprod U_j \to \mathcal{X}$ is a smooth surjective
morphism. Let $x : V \to \mathcal{X}$ be an object of $\mathcal{X}$. By
Sheaves on Stacks, Lemma
\ref{stacks-sheaves-lemma-surjective-flat-locally-finite-presentation}
there exists an \'etale covering $\{x_i \to x\}_{i \in I}$
such that each $x_i$ lifts to an object $u_i$ of $(\Sch/U)_\etale$.
This just means that $x_i$ lives over a scheme $V_i$, that
$\{V_i \to V\}$ is an \'etale covering, and that $x_i$ comes from
a morphism $u_i : V_i \to U$. Then
$x_i^*\mathcal{F} = u_i^*f^*\mathcal{F}$ is quasi-coherent.
This implies that $x^*\mathcal{F}$ on $(\Sch/V)_\etale$
is quasi-coherent, for example by
Modules on Sites, Lemma \ref{sites-modules-lemma-local-final-object}.
By Sheaves on Stacks, Lemma
\ref{stacks-sheaves-lemma-characterize-quasi-coherent-bis}
we see that $x^*\mathcal{F}$ is an fppf sheaf and since $x$
was arbitrary we see that $\mathcal{F}$ is a sheaf in the
fppf topology. Applying Sheaves on Stacks, Lemma
\ref{stacks-sheaves-lemma-characterize-quasi-coherent}
we see that $\mathcal{F}$ is quasi-coherent.

\medskip\noindent
Proof of (2). This is proved using exactly the same argument, which we fully
write out here. We may replace each of the algebraic stacks $\mathcal{X}_j$
by a scheme $U_j$ (using that any algebraic stack has a smooth covering by
a scheme and that flat and locally finite presented morphisms are preserved
by composition, see Morphisms of Stacks, Lemmas
\ref{stacks-morphisms-lemma-composition-flat} and
\ref{stacks-morphisms-lemma-composition-finite-presentation}).
The pullback of $\mathcal{F}$ to $(\Sch/U_j)_\etale$ is still
locally quasi-coherent, see
Sheaves on Stacks, Lemma \ref{stacks-sheaves-lemma-pullback-quasi-coherent}.
Then $f = \coprod f_j : U = \coprod U_j \to \mathcal{X}$ is a surjective,
flat, and locally finitely presented morphism. Let
$x : V \to \mathcal{X}$ be an object of $\mathcal{X}$. By
Sheaves on Stacks, Lemma
\ref{stacks-sheaves-lemma-surjective-flat-locally-finite-presentation}
there exists an fppf covering $\{x_i \to x\}_{i \in I}$
such that each $x_i$ lifts to an object $u_i$ of $(\Sch/U)_\etale$.
This just means that $x_i$ lives over a scheme $V_i$, that
$\{V_i \to V\}$ is an fppf covering, and that $x_i$ comes from
a morphism $u_i : V_i \to U$. Then
$x_i^*\mathcal{F} = u_i^*f^*\mathcal{F}$ is quasi-coherent.
This implies that $x^*\mathcal{F}$ on $(\Sch/V)_\etale$
is quasi-coherent, for example by
Modules on Sites, Lemma \ref{sites-modules-lemma-local-final-object}.
By Sheaves on Stacks, Lemma
\ref{stacks-sheaves-lemma-characterize-quasi-coherent}
we see that $\mathcal{F}$ is quasi-coherent.
\end{proof}

\end{document}
